%% appendix-applications -- joins in every tree class; the direct
%% reductions to plain entailment (flat bound witnesses, unipolar witnesses,
%% functional witnesses, the mixed sandwich) with proofs; and a covariant
%% entailment that holds for trees but fails in an extensible algebra.

\providecommand{\parity}{\operatorname{par}}

\section{Proofs for Constrained Types and Simplification}\label{app:ap}

The results of this appendix hold for $\BFT$ with any arity $n$ and variance
$c$, and for $\BF$ with $f$ covariant, in each tree class $\mathcal D$, with
the caveats stated. Terms are finite and built from variables, $\Bot$, $f$ and,
over $\BFT$, $\Top$; a term $s$ denotes the tree $\rho(s)$ under an assignment
$\rho$, and $s\le t$ means $\rho(s)\le_c\rho(t)$.

\begin{lemma}[Joins]\label{lem:ap-joins}
Any two trees $t,u$ have a least upper bound $t\vee u$ for $\le_c$, and $\Bot$ is
the least tree. If $t$ and $u$ are finite, or regular, then so is $t\vee u$.
If $f$ is covariant and $t$ and $u$ are free of $\Top$, then so is $t\vee u$.
\end{lemma}

\begin{proof}
The sets $L^0_t\cup L^0_u$ and $L^1_t\cap L^1_u$ are extension closed and
disjoint, since $L^0_t\cap L^1_t=L^0_u\cap L^1_u=\varnothing$, so by
\cref{thm:lf-representation}(1) they are the leaf pair of a tree $t\vee u$. By
\cref{thm:lf-representation}(2), a tree lies above $t$ and $u$ iff its first cone
contains $L^0_t\cup L^0_u$ and its second cone is contained in
$L^1_t\cap L^1_u$, so $t\vee u$ is the least upper bound. The leaf pair of
$\Bot$ is $(\varnothing,\Sigma^*)$ (\cref{thm:lf-representation}(3)), so $\Bot$
is least. A word outside both cones of $t\vee u$ lies outside $L^0_t\cup L^0_u$
and outside $L^1_t$ or $L^1_u$, so it is a position of $t$ or of $u$ labeled
$f$; by \cref{thm:lf-representation}(4) the join of finite trees is finite.
Unions and intersections of regular languages are regular. For covariant
$f$, $\nu_0$ is the identity, so a tree is free of $\Top$ iff its first cone is
empty, which the union preserves.
\end{proof}

\begin{proposition}[Flat bound witnesses]\label{prop:ap-flat}
Let
\[
 \psi=\psi_0\wedge\bigwedge_{(i,j)\in E}z_i\le z_j\wedge
 \bigwedge_{a\in L}s_a\le z_{i(a)}\wedge\bigwedge_{b\in U}z_{j(b)}\le t_b ,
\]
where $\psi_0$ is a conjunction of inequalities between terms over $X$, the
$s_a$ and $t_b$ are terms over $X$, and $E$ is a set of pairs of hidden
indices. Write $i\to^*j$ for reachability in the directed graph $E$. Then an
assignment $\rho:X\to\mathcal D$ extends to a $\mathcal D$-valued solution of
$\psi$ iff $\rho$ satisfies $\psi_0$ and $s_a\le t_b$ for all $a\in L$ and
$b\in U$ with $i(a)\to^*j(b)$. Hence $\varphi\models_{\mathcal D}\Exists\psi$
iff $\varphi$ entails, in $\mathcal D$, each literal of $\psi_0$ and each of
these inequalities.
\end{proposition}

\begin{proof}
If $\rho'$ is a solution of $\psi$ extending $\rho$, then $\rho$ satisfies
$\psi_0$, and for $i(a)\to^*j(b)$ the chain
$s_a\le z_{i(a)}\le\cdots\le z_{j(b)}\le t_b$ gives $s_a\le t_b$. Conversely,
put $\rho'(z_j)=\bigvee\{\rho(s_a):i(a)\to^*j\}$, with $\Bot$ for the empty
join; it lies in $\mathcal D$ by \cref{lem:ap-joins}. Each lower bound
$s_a\le z_{i(a)}$ holds since $i(a)\to^*i(a)$. For $(i,j)\in E$, every $a$
with $i(a)\to^*i$ has $i(a)\to^*j$, so $\rho'(z_i)\le_c\rho'(z_j)$. Each joinand
of $\rho'(z_{j(b)})$ is below $\rho(t_b)$ by hypothesis, hence so is the join.
The last statement follows because $\varphi\models_{\mathcal D}\Exists\psi$
says that every solution of $\varphi$ extends.
\end{proof}

With $q$ hidden variables, reachability costs $O(q^3)$, and at most
$\len L\cdot\len U$ plain entailments are added to those for $\psi_0$. The
\emph{occurrence sign} of a variable occurrence in an inequality $s\le t$
between terms is negative if the occurrence is in $s$ and positive if it is in
$t$, exchanged when the polarity $\parity(w)$ of its position $w$ in that term
is $1$.

\begin{proposition}[Unipolar witnesses]\label{prop:ap-unipolar}
Let $\psi$ be a conjunction of inequalities between terms over $X\cup Z$. Let
$Z_-$ be the hidden variables all of whose occurrences in $\psi$ are negative,
and $Z_+$ those, outside $Z_-$, all of whose occurrences are positive. Let
$\psi'$ be $\psi$ with every $z\in Z_-$ replaced by $\Bot$ and every $z\in Z_+$
replaced by $\Top$. Over $\BFT$, in each tree class,
$\varphi\models\Exists\psi$ iff $\varphi\models\exists\bar z'\,\psi'$, where
$\bar z'$ are the remaining hidden variables. Over $\BF$ the same holds with
$\Top$ replaced by the tree $\Omega$, whose constructor positions are all
words, over arbitrary and regular trees, and also over finite trees if
$Z_+=\varnothing$ or $n=0$.
\end{proposition}

\begin{proof}
By induction on terms, a term is monotone in a variable all of whose
occurrences in it have polarity $0$, and antitone in one all of whose
occurrences have polarity $1$: an occurrence at the position $iw$ of
$f(s_j)_j$ has polarity $c(i)\oplus\parity(w)$, and $f$ is monotone in a
covariant and antitone in a contravariant argument (\cref{sec:bg-trees}).
Hence increasing a variable whose occurrences in $s\le t$ are all positive
cannot falsify $s\le t$, and decreasing one whose occurrences are all negative
cannot either. Given a solution of $\psi$ extending a visible assignment,
change the values of the variables of $Z_-$ to $\Bot$ and those of $Z_+$ to
$\Top$, one at a time. Each change preserves every inequality, since
occurrence signs do not depend on the values of other variables, and $\Bot$
and $\Top$ lie in every tree class. The result solves $\psi'$ on the remaining
variables. Conversely, a solution of $\psi'$ gives a solution of $\psi$ with the
substituted values. Over $\BF$ the same argument uses $\Omega$, the greatest
tree over $\BF$, which is regular; for $n\ge1$ it is not finite, and for $n=0$
it is the finite tree $f()$.
\end{proof}

When every hidden variable is unipolar, $\psi'$ has no hidden variable, and the
entailment is a conjunction of plain entailments. Over $\BF$ the tree $\Omega$
is named in a plain query by a variable $o$ with $f(o,\ldots,o)\le o$, which
forces $\Omega$ over arbitrary and regular trees; for $n\ge1$, finite trees over
$\BF$ have no greatest element, and this elimination is not available there. The
occurrence signs are taken in the whole consequent: in a subsumption formula
the constraints $C_1$ and the context comparisons of $R_{12}$ count, not only
the result type.

\begin{proposition}[Functional witnesses]\label{prop:ap-functional}
Let $\psi=D\wedge R$, where $D$ consists of one equation $z=t_z$ for each
hidden variable $z$, with $t_z$ a term over $X\cup Z$, and $R$ is a conjunction
of inequalities between terms. The \emph{dependency graph} has an edge from $z$
to $z'$ when $z'$ occurs in $t_z$; the edge is \emph{guarded} unless $t_z=z'$.
Suppose that every cycle of the dependency graph contains a guarded edge. Over
arbitrary and over regular trees,
\[
 \varphi\models\Exists(D\wedge R)\iff\varphi\wedge D\models R,
\]
where the right side is a conjunction of plain entailments with $\bar z$
universally quantified. Over finite trees the same holds if the dependency
graph is acyclic. If it has a cycle, then $\varphi\models_{\mathrm{fin}}\Exists\psi$
iff $\varphi$ has no finite solution.
\end{proposition}

A guarded cycle forces infinite trees, so over finite trees the premise
$\varphi\wedge D$ of the plain entailments has no solution, and they would hold
vacuously; this is why the finite case needs its own dispatch.

\begin{proof}
\emph{Unique solutions.} An unguarded edge goes from $z$ to the variable $t_z$.
Since every cycle contains a guarded edge, the unguarded edges form an acyclic
graph in which each vertex has at most one outgoing edge. For a hidden $z$ let
$z^*$ be the end of its unguarded chain, so that $t_{z^*}$ is not a hidden
variable. Fix values $\rho(x)$ of the visible variables. For a term $s$ over
$X\cup Z$ and a word $w$ define a label $\lambda(s,w)$, possibly undefined: for
$s=x$ visible, the label of $\rho(x)$ at $w$; for a constant, its label at
$w=\eps$; for $s=f(s_j)_j$, the label $f$ at $\eps$ and $\lambda(s_i,w')$ at
$w=iw'$; and for $s=z$ hidden, $\lambda(t_{z^*},w)$. Each step either ends,
shortens $w$, or passes from a hidden variable to a term that is not a hidden
variable, so the recursion terminates. By induction on the length of $w$, the
labels $\lambda(s,\cdot)$ define trees $T(s)$ with $T(f(s_j)_j)=f(T(s_j))_j$ and
$T(z)=T(t_z)$, so the values $T(z)$ solve $D$, and every solution of $D$ with
these visible values agrees with $T$ at every position, by the same induction.
Every subtree of $T(z)$ is $T(s)$ for a subterm $s$ of some $t_{z'}$, a subtree of
some $\rho(x)$, or a constant, so $T(z)$ is regular when every $\rho(x)$ is. If
the dependency graph is acyclic, then $T(z)$ is the value of the finite term
obtained from $z$ by substituting definitions finitely often, which is finite
when every $\rho(x)$ is. If the graph has a cycle
$z_1\to z_2\to\cdots\to z_k=z_1$, let $w_i$ be a position of $z_{i+1}$ in
$t_{z_i}$. In every solution the subtree of $z_i$ at $w_i$ is the value of
$z_{i+1}$, since the path to an occurrence in a term passes through
constructors. Some $w_i$ is nonempty, so $z_1$ is its own subtree at the
nonempty position $w=w_1\cdots w_{k-1}$, has all positions $w^j$, and is
infinite.

\emph{The reduction.} By the above, each assignment $\rho$ of $X$ has exactly
one extension to a solution of $D$, and it lies in the tree class of $\rho$ in
the stated cases. So $\rho$ extends to a solution of $D\wedge R$ iff this unique
extension satisfies $R$, and the solutions of $\varphi\wedge D$ are exactly the
solutions of $\varphi$ with their unique extensions. This proves the
equivalence over arbitrary and regular trees, and over finite trees when the
graph is acyclic. If it has a cycle, $D$ has no finite solution for any visible
values, so finite entailment holds iff $\varphi$ has no finite solution.
\end{proof}

\begin{proposition}[Mixed sandwich]\label{prop:ap-sandwich}
For arrow types, $\varphi\models_{\mathcal D}\exists z\,[\,x\le(z\to z)\ \wedge\
(z\to z)\le y\,]$ iff $\varphi\wedge x=\Top$ and $\varphi\wedge y=\Bot$ have no
solution in $\mathcal D$ and the following plain entailments hold in
$\mathcal D$, where $a,b,c',d'$ are fresh visible variables:
\[
\begin{aligned}
 &\varphi\wedge x=(a\to b)\wedge y=\Top\models b\le a,\qquad
 \varphi\wedge x=\Bot\wedge y=(c'\to d')\models c'\le d',\\
 &\varphi\wedge x=(a\to b)\wedge y=(c'\to d')\models
 b\le a\wedge b\le d'\wedge c'\le a\wedge c'\le d' .
\end{aligned}
\]
\end{proposition}

The hidden variable $z$ occurs at both polarities of $z\to z$, so the
consequent has no parity assignment and $z$ is not unipolar; still, two
satisfiability tests and six plain entailments decide the instance.

\begin{proof}
Let $\rho$ be a solution of $\varphi$. If $\rho(x)=\Top$ or $\rho(y)=\Bot$, no
$z$ exists, because $z\to z$ has the root $f$. Otherwise $\rho(x)$ is $\Bot$ or
an arrow $a\to b$, and $\rho(y)$ is $\Top$ or an arrow $c'\to d'$. By the
order on arrows, $a\to b\le_c z\to z$ iff $z\le_c a$ and $b\le_c z$, and
$z\to z\le_c c'\to d'$ iff $c'\le_c z$ and $z\le_c d'$, while the comparisons with
$\Bot$ and $\Top$ always hold. So the witnesses form the set of $z$ above the
lower bounds and below the upper bounds of the four cases, and by
\cref{prop:ap-flat}, with one hidden variable, it is nonempty iff every lower
bound lies below every upper bound: no condition, $b\le a$, $c'\le d'$, or the
four comparisons. The case premises are mutually exclusive and cover the
solutions of $\varphi$ with $\rho(x)\ne\Top$ and $\rho(y)\ne\Bot$, and each
solution of $\varphi$ extends in exactly one way to the fresh variables of its
case, which then denote the children. Hence the case conditions hold in every
solution of $\varphi$ iff the plain entailments hold, and the two
satisfiability tests say that no solution has $\rho(x)=\Top$ or $\rho(y)=\Bot$.
\end{proof}

\subsection{Closed Trees and Extensible Algebras}\label{app:ap-closed}

Let $f$ be binary and covariant over $\BFT$, and
\[
 \Phi=\{x\le f(a,e),\ a\le y,\ f(m,e)\le y,\ f(e,e)\le m,\ \Top\le e\}.
\]
Then $\Phi\models x\le y$ over trees. Indeed $e=\Top$. If $y=\Top$ there is
nothing to prove. Otherwise $f(m,\Top)\le y$ forces $y=f(y_0,\Top)$ with
$m\le y_0$, and $a\le y$ forces $a\ne\Top$, so
$a\le f(\Top,\Top)\le m\le y_0$. Hence $x\le f(a,\Top)\le f(y_0,\Top)=y$. Thus
deleting $x\le y$ from $\Phi\wedge x\le y$ is an exact simplification, decided
by one plain entailment, or by \cref{prop:ap-flat} with the witness $z=x$ for
$\exists z\,(x\le z\le y)$.

Now consider the functions from $\{0,1\}^*$ to the subsets of
$\{\mathsf a,\mathsf b\}$, ordered pointwise by inclusion, with $\Bot$ the
constant $\varnothing$, $\Top$ the constant $\{\mathsf a,\mathsf b\}$, and
$f(p,p')$ the function with the value $\{\mathsf a\}$ at $\eps$, $p(w)$ at $0w$
and $p'(w)$ at $1w$. Pointwise unions and intersections make this a
distributive lattice, and $f$ preserves binary joins and meets. A tree $t$
embeds as the function whose value at $w$ contains $\mathsf a$ iff
$w\notin L^1_t$ and contains $\mathsf b$ iff $w\in L^0_t$: the value is
$\{\mathsf a\}$ at the positions labeled $f$, $\{\mathsf a,\mathsf b\}$ below a
leaf $\Top$, and $\varnothing$ below a leaf $\Bot$. By
\cref{thm:lf-representation}(2,3) this embedding preserves and reflects the
order and commutes with $f$, $\Bot$ and $\Top$. Let $q$ have the value
$\{\mathsf b\}$ at $\eps$ and $\varnothing$ elsewhere, and put $e=\Top$,
$m=f(\Top,\Top)$, $a=q$, $y=q\vee f(m,\Top)$ and $x=f(q,\Top)$. All literals of
$\Phi$ hold, but $x(0)=q(\eps)=\{\mathsf b\}$ and
$y(0)=q(0)\cup m(\eps)=\{\mathsf a\}$, so $x\not\le y$. Hence a procedure that is
sound for this extensible algebra cannot certify the simplification; this
separates the two semantics and does not indicate an error in a procedure
specified for the extensible one.
